\cotefile{157-4}
\scanpage{394}\mspage{118}
fondamentaux de $\mathrm{Cat}$) montre à
nouveau que si $W$ \inserted{loc. fond.} de $\mathrm{Cat}$ satisfait
à $\mathrm{Loc}(6'\mathrm{a},8,9)$, alors $\mathrm{HOT}_{W}$ est
un vrai dérivateur. On a vu de
plus dans 7.6 (page 111) que l'on
a \inserted{alors} $W=W^{\ast}_{\mathrm{HOT}_{W}}$, donc il s'ensuit
\deleted{que}
% slanted
\margin{Mais c'était sous l'hypothèse $\mathrm{Loc}(3\ \mathrm{bis})$ ! \unsure{Gare au|Erreur} cercle vicieux !}

\uline{Corollaire} 7.7.2 : \deleted{Les axiomes $\mathrm{Loc}(6\mathrm{a},8,9)$} \inserted{On a l'implication}
\[
\mathrm{Loc}(6'\mathrm{a},8,9)\Longrightarrow\mathrm{Loc}(1\ \mathrm{ter},\deleted{3\ \mathrm{bis}}).
\]
% slanted
\margin{pas prouvé, probablement faux}

Notons aussi, pour mémoire :

\uline{Corollaire} 7.7.3. Soit $(M,W)$ une
catégorie de modèles, associée à
une catégorie à cofibrations. Alors
sur le domaine $\mathrm{Dia}$ formé des
ens ordonnés finis, le prédérivateur
\inserted{$D={}$}$D_{M,W}$ satisfait aux axiomes
$\mathrm{Der}(1,2,3,4^{\unsure{\alpha|\prime|\wedge|\mathrm{a}}},5^{\alpha},\unsure{6|5|0}^{\alpha})$, et dualement. Donc
si $W$ est associée à la fois à une
structure de cofibrations et de fibrations,
$D$ est un dérivateur sur $\mathrm{Dia}$.
% archive sheet number: 393
